
This paper presents a systematic investigation of when uncertainty-based hallucination detector calibration transfers across benchmarks. The experiments found that QA benchmarks cluster into two families with silhouette score 0.8245. Within-cluster threshold transfer showed 3.2\% mean AUROC degradation while cross-cluster transfer showed 22.3\% degradation---an approximately $7\times$ difference in the tested pairs.

These findings suggest that JS-divergence clustering may provide a criterion for predicting transfer success: check distribution similarity before deployment and recalibrate when distributions differ substantially. However, these results are based on proof-of-concept experiments with limited samples and benchmark pairs.

Future work should extend this framework to larger models, additional benchmark families, full-scale validation across all benchmark pairs, and API-only models requiring alternative uncertainty methods.

